\documentclass[a4paper,11pt]{article}
\usepackage[T1]{fontenc}\usepackage[utf8]{inputenc}\usepackage{lmodern}\usepackage{microtype}
\usepackage[a4paper,margin=1in,headheight=15pt,headsep=14pt]{geometry}
\usepackage{amsmath,amssymb}\usepackage{graphicx}\usepackage{booktabs}
\usepackage[table,svgnames]{xcolor}\usepackage[most]{tcolorbox}\tcbuselibrary{skins,breakable}
\usepackage{tikz}\usepackage{pifont}\usepackage{bbding}\usepackage{enumitem}
\setlist{leftmargin=1.5em,topsep=4pt,itemsep=3pt}\usepackage{parskip}
\usepackage{fancyhdr}\usepackage{titlesec}\usepackage[hidelinks]{hyperref}\usepackage{xurl}
\hypersetup{breaklinks=true}
\definecolor{navy}{HTML}{21355E}\definecolor{navyrule}{HTML}{21355E}
\definecolor{inktext}{HTML}{1A202C}\definecolor{mutedtext}{HTML}{6B7280}
\definecolor{intuFrame}{HTML}{2C5AA0}\definecolor{intuFill}{HTML}{F4F7FC}
\definecolor{everyFrame}{HTML}{8A5A1E}\definecolor{everyFill}{HTML}{FBF1E3}
\definecolor{workFrame}{HTML}{2E7D52}\definecolor{workFill}{HTML}{F1F8F3}
\definecolor{watchFrame}{HTML}{B23A48}\definecolor{watchFill}{HTML}{FBF1F2}
\definecolor{keyFrame}{HTML}{6A4C93}\definecolor{keyFill}{HTML}{F4F0F9}
\color{inktext}\hypersetup{linkcolor=navy,urlcolor=navy,citecolor=navy}
\tcbset{housecallout/.style 2 args={enhanced,breakable,colback=#2,colframe=#1,boxrule=0.6pt,arc=2.4pt,outer arc=2.4pt,left=10pt,right=10pt,top=9pt,bottom=9pt,before skip=11pt,after skip=11pt,fontupper=\normalsize,coltitle=white,fonttitle=\bfseries\sffamily\small,attach boxed title to top left={xshift=6mm,yshift=-3mm},boxed title style={colback=#1,colframe=#1,rounded corners,arc=3pt,boxrule=0pt,left=6pt,right=6pt,top=2pt,bottom=2pt}}}
\newtcolorbox{intuition}{housecallout={intuFrame}{intuFill},title={\raisebox{-0.05ex}{\Pointinghand}\,\ The intuition}}
\newtcolorbox{everyday}{housecallout={everyFrame}{everyFill},title={\raisebox{-0.05ex}{$\bigstar$}\,\ Everyday picture}}
\newtcolorbox{worked}[1]{housecallout={workFrame}{workFill},title={\raisebox{-0.05ex}{\Pencil}\,\ Worked example: #1}}
\newtcolorbox{watchout}{housecallout={watchFrame}{watchFill},title={\raisebox{-0.05ex}{\ding{55}}\,\ Watch out}}
\newtcolorbox{keytake}{housecallout={keyFrame}{keyFill},title={\raisebox{-0.05ex}{\ding{51}}\,\ Key takeaway}}
\newcommand{\CourseShort}{SEML Companion}\newcommand{\SessionNo}{6}\newcommand{\LectureTitle}{EDA \& ML Design Patterns}
\pagestyle{fancy}\fancyhf{}
\fancyhead[L]{\footnotesize\color{mutedtext}\textsf{\CourseShort}}
\fancyhead[R]{\footnotesize\color{mutedtext}\textsf{Session \SessionNo\ \textperiodcentered\ \LectureTitle}}
\fancyfoot[C]{\footnotesize\color{mutedtext}\thepage}
\renewcommand{\headrulewidth}{0.4pt}\renewcommand{\footrulewidth}{0pt}
\renewcommand{\headrule}{\hbox to\headwidth{\color{navyrule!55}\leaders\hrule height \headrulewidth\hfill}}
\titleformat{\section}{\sffamily\Large\bfseries\color{navy}}{\thesection}{0.8em}{}[{\vspace{2pt}\color{navy!70}\titlerule[0.8pt]}]
\titleformat{\subsection}{\sffamily\large\bfseries\color{navy!92}}{\thesubsection}{0.7em}{}
\titlespacing*{\section}{0pt}{16pt}{8pt}\titlespacing*{\subsection}{0pt}{12pt}{4pt}
\newcommand{\titlebanner}[4]{\begin{tcolorbox}[enhanced,breakable=false,colback=navy,colframe=navy,boxrule=0pt,arc=3pt,outer arc=3pt,left=16pt,right=16pt,top=16pt,bottom=16pt,before skip=2pt,after skip=14pt,width=\linewidth]{\sffamily\footnotesize\bfseries\color{white!85}\MakeUppercase{#1}}\par\vspace{6pt}{\sffamily\bfseries\fontsize{30}{34}\selectfont\color{white}#2\par}\vspace{5pt}{\sffamily\large\color{white!88}#3\par}\vspace{9pt}{\color{white!55}\rule{\linewidth}{0.4pt}}\par\vspace{6pt}{\sffamily\footnotesize\color{white!82}#4\par}\end{tcolorbox}}
\newcommand{\vocab}[1]{\textbf{\color{navy}#1}}
\begin{document}
\thispagestyle{fancy}
\titlebanner{AIMLCZG546 · Session 6 · Companion Reader}{Event-Driven Architecture \& ML Design Patterns}{Model Registry, Batch Serving, and Real-Time Inference Explained}{Session 6 · Read alongside the lecture slides · Every pattern and concept worked through in full}

\noindent\textbf{How to use this companion.} \textcolor{intuFrame}{\textbf{Blue}} = intuition. \textcolor{everyFrame}{\textbf{Amber}} = everyday analogy. \textcolor{workFrame}{\textbf{Green}} = worked scenario. \textcolor{watchFrame}{\textbf{Red}} = pitfall. \textcolor{keyFrame}{\textbf{Purple}} = one line to remember.

%==============================================================
\section{Event-Driven Architecture}
%==============================================================

\vocab{Event-Driven Architecture (EDA)} is an architectural style in which components communicate by producing and consuming \emph{events} rather than calling each other directly.

\textbf{Context:} A complex system handles asynchronous events from user actions, external systems, or internal processes. Services must react to these events without being tightly coupled to each other.

\textbf{Problem:} If Service A directly calls Service B, they are coupled. If Service B is slow or down, Service A is blocked. If you add Service C (also interested in the same events), you must modify Service A again.

\textbf{Solution:} Introduce a \vocab{message broker} between producers and consumers.

\begin{intuition}
Imagine a school PA system. The principal (publisher) makes one announcement. Every classroom (subscriber) hears it simultaneously. The principal does not know or care which classrooms exist or how they react. Adding a new classroom does not require the principal to change anything.
\end{intuition}

\subsection{Publish-Subscribe Model}

\begin{itemize}
  \item \vocab{Publishers} (Producers) --- generate events and publish them to the broker without knowing who will consume them.
  \item \vocab{Message Broker} --- stores, routes, and delivers events. Examples: Apache Kafka, RabbitMQ, Amazon SQS, Azure Service Bus.
  \item \vocab{Subscribers} (Consumers) --- declare interest in specific event types; the broker delivers matching events.
\end{itemize}

\begin{worked}{Fraud detection with EDA}
A transaction occurs at 14:32:01 for customer C491, amount \$1,200, merchant: ``Electronics Ltd, Lagos.'' \\[4pt]
\textbf{Step 1 --- Payment Service} (Publisher) publishes event to Kafka topic \texttt{transactions.created}: \\
\texttt{\{"id":"txn-9921","customer":"C491","amount":1200,"merchant":"Electronics Ltd","country":"NG"\}} \\[4pt]
\textbf{Step 2 --- Fraud Detection Service} (Subscriber) consumes the event. Runs XGBoost model. Fraud probability = 0.94. Publishes to topic \texttt{fraud.detected}: \\
\texttt{\{"txn\_id":"txn-9921","fraud\_prob":0.94,"action":"BLOCK"\}} \\[4pt]
\textbf{Step 3 --- Notification Service} (Subscriber to \texttt{transactions.created}) sends SMS: ``A transaction of \$1,200 in Nigeria was attempted. Was this you?'' \\[4pt]
\textbf{Step 4 --- Analytics Service} (Subscriber to \texttt{transactions.created}) updates the fraud dashboard. \\[4pt]
\textbf{Why EDA wins here:} The Payment Service published one event. Three services reacted independently. Adding a fourth (e.g., a compliance logging service) requires zero changes to existing services --- just subscribe to the topic.
\end{worked}

\noindent\textbf{Where this shows up in ML/AI.} ML model retraining pipelines are naturally event-driven: ``data drift detected'' event triggers retraining; ``new labelled batch available'' event triggers evaluation; ``model deployed'' event triggers A/B test initialisation.

%==============================================================
\section{Model Registry Pattern}
%==============================================================

\subsection{The Registry Pattern (Software Engineering)}

A \vocab{registry} is a central hub that dynamically stores, locates, and retrieves class instances or services. It acts as a directory: ``where is the thing I need?'' Combined with \vocab{version control}, it becomes a model registry.

\subsection{What a Model Registry Stores}

\begin{center}
\resizebox{\linewidth}{!}{%
\begin{tabular}{ll}
\toprule
\textbf{Git (code version control)} & \textbf{ML Model Registry} \\
\midrule
Source code files (.py, .java) & Trained model files (.pkl, .pt, .onnx) \\
Documentation (.md, .rst) & Model metrics (accuracy, F1, AUC by dataset version) \\
Configuration files & Hyperparameters and training configuration \\
Commit messages & Experiment run description and notes \\
Branch/tag names & Model aliases (staging, production, champion) \\
\multicolumn{1}{l}{---} & Dataset version and data lineage \\
\multicolumn{1}{l}{---} & Training environment (Python version, package versions) \\
\bottomrule
\end{tabular}}
\end{center}

\subsection{MLflow Registry in Practice}

MLflow is the most widely used open-source ML registry. Its four core capabilities:

\begin{enumerate}
  \item \textbf{Version Control} --- every trained model is saved as a versioned artefact. Multiple versions co-exist; you can compare, roll back, or deploy any version.
  \item \textbf{Lineage and Traceability} --- each model version links to the exact training run (data version, hyperparameters, metrics, code commit).
  \item \textbf{Production Workflows} --- \vocab{aliases} and \vocab{tags} control the promotion pipeline: a model moves from \texttt{@staging} to \texttt{@production} only after evaluation passes.
  \item \textbf{Governance} --- access control, audit log, metadata for regulatory compliance.
\end{enumerate}

\begin{worked}{Churn model lifecycle in MLflow}
\textbf{Day 1:} Data scientist trains v1 (logistic regression). Logs to MLflow: accuracy=0.76, AUC=0.79. Registers as version 1. Tag: \texttt{status=experimental}. \\[4pt]
\textbf{Day 5:} Trains v2 (XGBoost, new features). Logs: accuracy=0.84, AUC=0.88. Registers as version 2. Sets alias: \texttt{@staging}. \\[4pt]
\textbf{Day 7:} QA team runs offline evaluation on holdout set. Version 2 passes. Sets alias: \texttt{@production}. Version 1 alias cleared. \\[4pt]
\textbf{Day 30:} New team trains v3 (LightGBM, fresh data). Registers as version 3. Sets alias: \texttt{@staging}. Begins A/B test in production (5\,\% traffic to v3). \\[4pt]
\textbf{Day 37:} A/B test complete. Version 3 reduces churn by 8\,\% vs version 2. Sets \texttt{@production} to version 3. Version 2 demoted. \\[4pt]
\textbf{Rollback scenario:} Version 3 causes a p95 latency spike in production. Team runs: \texttt{mlflow.set\_registered\_model\_alias("churn-model","production","2")}. Traffic immediately routed back to version 2. Total rollback time: 45 seconds.
\end{worked}

%==============================================================
\section{Batch vs Real-Time Serving}
%==============================================================

Two fundamentally different patterns for deploying ML models in production. The choice depends on latency requirements, throughput, cost, and business context.

\subsection{Batch Prediction (Offline)}

The model runs periodically over a large dataset and saves predictions for later consumption.

\begin{everyday}
A newspaper is printed once per day with tomorrow's weather forecast. Readers get the information with a delay, but the cost of printing is amortised over millions of readers. You can't ask the newspaper to predict tomorrow's weather right now, at this exact second.
\end{everyday}

\textbf{When to use Batch:}
\begin{itemize}
  \item Customer segmentation: group users by weekly behaviour → email campaign.
  \item Credit risk scoring at end of business day.
  \item Product demand forecasting for warehouse stocking (run nightly).
  \item Recommendation emails: score all users at 2am, send emails at 9am.
\end{itemize}

\subsection{Online Prediction (Real-Time)}

The model is deployed as a persistent REST API and returns predictions in milliseconds, on demand.

\begin{everyday}
Weather radar: the moment you open the app, it queries the satellite and shows you what is happening \emph{right now}. Instant response is the point; a 24-hour delay is useless.
\end{everyday}

\textbf{When to use Online:}
\begin{itemize}
  \item Fraud detection: must block a transaction \emph{before} it is authorised.
  \item Dynamic pricing: ride-share surge pricing changes every few seconds.
  \item Spam filtering: classify before the email lands in the inbox.
  \item Chatbot responses: user is waiting.
\end{itemize}

\subsection{Comparison}

\begin{center}
\resizebox{\linewidth}{!}{%
\begin{tabular}{lll}
\toprule
\textbf{Property} & \textbf{Batch} & \textbf{Online} \\
\midrule
Latency & Minutes to hours & Milliseconds to seconds \\
Throughput & Very high (process millions per run) & Lower (request-by-request) \\
Cost & Low (run only when needed) & Higher (server always running) \\
Data freshness & Stale (predictions age between runs) & Real-time (uses current data) \\
Trigger & Scheduled (cron, Airflow) & Per request \\
Infrastructure & Spark, Airflow, BigQuery & FastAPI, Kubernetes, Redis \\
\bottomrule
\end{tabular}}
\end{center}

\begin{worked}{Fraud detection: choosing the pattern}
Transaction arrives: customer paying \$150 at a POS terminal. Is this fraud? \\[4pt]
\textbf{Decision tree:} \\
Q1: Does the user need a response in $<$1 second? \textbf{Yes} (card is physically waiting). \\
Q2: Does the decision block a business transaction? \textbf{Yes} (payment must be approved or declined immediately). \\
$\Rightarrow$ \textbf{Online serving required.} \\[4pt]
\textbf{Contrast --- promotional email targeting:} \\
Q1: Does the user need a response in $<$1 second? \textbf{No} (email goes out tomorrow morning). \\
Q2: Does the decision block a business transaction? \textbf{No} (email is informational). \\
$\Rightarrow$ \textbf{Batch serving is sufficient and cheaper.}
\end{worked}

\begin{keytake}
Batch if the user can wait and data volume is large. Online if the decision must be made before the user moves on or a transaction completes.
\end{keytake}

\end{document}
